\section{Conclusion}

Tensor networks having been proposed as promising language models, we first apply TT decomposition to real-world language modeling datasets and name the framework TTLM. We propose two variants: TTLM-Large and TTLM-Tiny, and show that they outperform Vanilla-RNNs with low-scale hidden units. The presentation of the experimental results is an advancement for exploring tensor networks in machine learning. Meanwhile, we demonstrate that Second-order RNNs, RACs, and MI-RNNs are special implementations of TTLM. 

A limitation of this study is that it shall examine the influence of different normalization functions. In future research, if appropriate mathematical tools and benchmarks are available, we plan to investigate the long-range correlation modeling capability of TTLM in natural language, which is believed to be one of the core features in TTLM. 
% \benyou{no benchmark}











% In the unusual situation where you want a paper to appear in the
% references without citing it in the main text, use \nocite
% \nocite{langley00}

\bibliography{icml2023}
\bibliographystyle{icml2023}


%%%%%%%%%%%%%%%%%%%%%%%%%%%%%%%%%%%%%%%%%%%%%%%%%%%%%%%%%%%%%%%%%%%%%%%%%%%%%%%
%%%%%%%%%%%%%%%%%%%%%%%%%%%%%%%%%%%%%%%%%%%%%%%%%%%%%%%%%%%%%%%%%%%%%%%%%%%%%%%
% APPENDIX
%%%%%%%%%%%%%%%%%%%%%%%%%%%%%%%%%%%%%%%%%%%%%%%%%%%%%%%%%%%%%%%%%%%%%%%%%%%%%%%
%%%%%%%%%%%%%%%%%%%%%%%%%%%%%%%%%%%%%%%%%%%%%%%%%%%%%%%%%%%%%%%%%%%%%%%%%%%%%%%
